\section{Organisation fédérale et prérogatives}

\begin{frame}{La FFESSM}
  \begin{itemize}
    \item \textbf{Fédération Française d’Études et de Sports Sous-Marins}.
    \item Seule fédération reconnue par l’État pour :
          \begin{itemize}
            \item définir les \textbf{règles techniques et administratives},
            \item organiser les \textbf{compétitions} subaquatiques.
          \end{itemize}
    \item Regroupe les clubs associatifs et leurs licenciés.
    \item La licence permet : formations, brevets, assurance plongée, etc.
  \end{itemize}
\end{frame}

\begin{frame}{La CMAS}
  \begin{itemize}
    \item \textbf{Confédération Mondiale des Activités Subaquatiques}.
    \item Réunit des fédérations de nombreux pays.
    \item Objectifs : développer la connaissance et la protection du monde subaquatique,
          favoriser les sports subaquatiques.
    \item Votre Niveau 1 FFESSM correspond à une \textbf{étoile CMAS},
          reconnu dans de nombreux pays.
  \end{itemize}
\end{frame}

\begin{frame}{Prérogatives du plongeur Niveau 1}
  \begin{itemize}
    \item Accès à l’\textbf{espace médian} : jusqu’à \textbf{20 m}.
    \item Toujours \textbf{accompagné} par un guide au minimum Niveau 4.
    \item Aucune compétence d’\textbf{autonomie}, sauf pour ceux qui sont \textbf{PA12}.
  \end{itemize}

  \begin{block}{Conditions administratives}
    \begin{itemize}
      \item Certificat médical de non contre-indication (moins d’un an).
      \item Licence FFESSM en cours de validité.
      \item Passeport ou carte CMAS.
      \item Carnet de plongée.
      \item Autorisation parentale si mineur.
    \end{itemize}
  \end{block}
\end{frame}
